
\subsection{tipi}

V programski jezik sem vključila naslednje tipe
\begin{itemize}
    \item{\cd{LInt}} Linearno celo število. Za razliko od običajnega tipa \cd{Int}, se \cd{LInt} obnaša linearno. Torej po izvedbi kode
    \begin{lstlisting}[language=Java]
let a = 0
let b = a
    \end{lstlisting}
    je vir \cd{a : LInt} porabljen. Kaj pa neimenovane spremenljivke? Če je \cd{1 : LInt}, kako lahko računamo \cd{1 + 1}? Tu \cd{1} in \cd{1} nista dve pojavitvi istega elementa, temveč stva dva različna elementa tipa \cd{LInt}. Postavimo naslednje dva pravila
    \begin{center}
        \begin{bprooftree}
            \AXC{}
            \UIC{\cd{LInt} Type}
        \end{bprooftree}
        \qquad
        \begin{bprooftree}
            \AXC{}
            \UIC{\(n : \cd{LInt}\)}
        \end{bprooftree}
    \end{center}
    Torej tako kot lahko enoto \(\unit\) s pravilom vpeljave za enoto ``ustvarimo'' iz ničesar, lahko ``ustvarimo'' elemente tipa \cd{LInt}.

    \item{\cd{A -> B}} Linearni funkcijski tip. Ta tip predstavlja \(A \multimap B\). Ko apliciramo nek element \cd{a : A} na funkcijo \cd{f : A -> B} se vir \cd{a} uporabi. Poglejmo si primer v kodi
    \begin{lstlisting}
        let f x = x + 1
        let a = 3
        let b = f x
    \end{lstlisting}
\end{itemize}


\subsection{pravila v Line-u}


Pravila vpeljave tipov:
%
\begin{center}
    \begin{bprooftree}
        \AXC{}
        \UIC{\cd{LInt} : Type}
    \end{bprooftree} 
    \qquad
    \begin{bprooftree}
        \AXC{}
        \UIC{\cd{A -> B} : Type}
    \end{bprooftree}   
\end{center}

Za vsak \(n \in \mathbb{Z}\) imamo pravilo:
%
\begin{center}
        \begin{bprooftree}
            \AXC{}
            \UIC{\(\vdash n\) : \cd{LInt}}
        \end{bprooftree}
\end{center}




\begin{center}
    \qquad
    \begin{bprooftree}
        \AXC{\cd{let f x = b}???????????}
        \UIC{}
    \end{bprooftree}
    \qquad
    \begin{bprooftree}
        \AXC{\cd{f:A\(\multimap\)B}}
        \AXC{\cd{a:A}}
        \BIC{\cd{let f a = b???????????
f:A⊸B f a :B}}
    \end{bprooftree}
\end{center}